\ifx\danslelivre\undefined
\documentclass[10pt]{../fichiers-configuration-latex/smfart}
\usepackage{ifthen}
\newboolean{english}
\newboolean{francais}
\setboolean{english}{false}
\setboolean{francais}{true}
\input{../fichiers-configuration-latex/paquets}
\selectlanguage{francais}
\input{../fichiers-configuration-latex/lemme-proposition-theoreme}
\input{../fichiers-configuration-latex/numerotation}
\input{../fichiers-configuration-latex/macros}

\externaldocument{../00-introduction/introduction}
\externaldocument{../01-quasi-excellents/excellent}
\externaldocument{../02-topologies/topologies}
\externaldocument{../03-approximation/approximation}
\externaldocument{../04-Cohen-Gabber/Cohen-Gabber}
\externaldocument{../05-algebrisation-partielle/algebrisation}
\externaldocument{../06-log-lissite/log}
\externaldocument{../07-uniformisation-locale/uniformisation-faible}
%\externaldocument{../08-modification-theorem/modification}
\externaldocument{../12-descente-cohomologique-orientee/descente-orientee}
\externaldocument{../13-theoreme-de-finitude/finitude}
\externaldocument{../15-Lefschetz-affine/Lefschetz-affine}
\externaldocument{../16-classes-de-Chern/gysin}
\externaldocument{../17-complexes-dualisants/dualite}
\externaldocument{../19-exemples-contre-exemples/contre-exemples}
\externaldocument{../20-rigidite/rigidite}
\externaldocument{../21-finitude-non-abelienne/finitude-non-abelienne}


\begin{document}

\begin{center}
XIV. Fonctions de dimension\\
Vincent Pilloni et Benoît Stroh
\end{center}

\else
\setboolean{english}{false}
\setboolean{francais}{true}
\selectlanguage{francais}
\chapter{Fonctions de dimension}
\addtocontents{toc}{par Vincent Pilloni et Benoît Stroh}
\begin{center}
Vincent Pilloni et Benoît Stroh
\end{center}
\fi


Nous définissons la notion de fonction de dimension sur un schéma $X$ et nous montrons l'existence
de telles fonctions localement pour la topologie étale si $X$ est quasi-excellent.


\section{Universelle caténarité des schémas henséliens}

Dans cette partie, nous rappelons les notions de \emph{caténarité} et d'\emph{universelle
caténarité}.
Le lecteur pourra consulter l'exposé~I pour plus de détails.

\subsection{Schémas universellement caténaires}


Soient $S$ un espace topologique et $X\subset Y$ des fermés irréductibles de $S$. Notons
$\codim(X,Y)$ la borne supérieure de l'ensemble des longueurs des chaînes strictement croissantes
de fermés irréductibles $X\varsubsetneq Z \subset Y$. Si $S$ est un schéma,  $X$ et $Y$ des
sous-schémas fermés intègres et $x$ le point générique de $X$, on a
$$\codim(X,Y)\:=\:\dim(\OO_{Y,x})\: .$$

\begin{definition2}[\egazero{iv}{14.3.2}] \label{14-definition-catenaire} Un schéma $S$ est
\emph{caténaire} s'il est nœthérien et si pour toute chaîne $X\subset Y\subset Z$ de fermés
irréductibles de $S$, on a $$\codim(X,Z)=\codim(Y,Z)+\codim(X,Y).$$ Un schéma $S$ est
\emph{universellement caténaire} si tout schéma de type fini sur $S$ est caténaire.
\end{definition2}

La notion de caténarité est stable par localisation et par restriction à des sous-schémas fermés.
Ainsi, $S$ est universellement caténaire si et seulement si pour tout entier $n\geq 0$, le schéma
$\Aff^n_S$ est caténaire.

\begin{lemme2} \label{14-lemme-CohenMac-implique-univ-cat}  Un schéma de Cohen-Macaulay est
universellement caténaire.  \end{lemme2}

\begin{démo} Si $S$ est Cohen-Macaulay, il est caténaire d'après \cite{CA@Matsumura}~16.B. Comme
pour tout $n\geq 0$, le schéma $\Aff^n_S$ reste Cohen-Macaulay, le schéma $S$ est bien
universellement caténaire.  \end{démo}

\begin{exemple2}\label{14-exemple-regulier-implique-univ-cat}  Tout schéma régulier est
universellement caténaire car Cohen-Macaulay. En particulier, le spectre d'un corps, un trait et le
spectre d'une algèbre de séries formelles sur un corps ou sur un anneau de valuation discrète sont
universellement caténaires. Tout schéma de type fini sur un schéma régulier est universellement
caténaire.  \end{exemple2}

\begin{proposition2}[\cite{CA@Matsumura}~28.P] \label{14-prop-complet-univ-cat} Un schéma local
complet nœthérien est universellement caténaire.  \end{proposition2}

\begin{démo} Le théorème de structure de Cohen~\egazero{iv}{19.8.8} permet d'écrire tout schéma
local complet nœthérien comme fermé dans le spectre d'une algèbre de séries formelles sur un anneau
de Cohen. L'universelle caténarité résulte de l'exemple précédent et de la stabilité de cette notion
par passage à un fermé.  \end{démo}

\subsection{Un théorème de Ratliff}

On dit qu'un schéma nœthérien est \emph{équidimensionnel} si toutes ses composantes irréductibles
ont même dimension (finie). Soit $S$ un schéma local nœthérien. On note $\widehat{S}$ le spectre du
complété de l'anneau de $S$ en son idéal maximal.

\begin{definition2} \label{14-defi-formellement-catenaire} Le schéma local $S$ est
\emph{formellement équidimensionnel} si  $\widehat{S}$ est équidimensionnel. Il est
\emph{formellement caténaire} si pour tout $s\in S$, l'adhérence $\overline{\{s\}}$ est formellement
équidimensionnelle.  \end{definition2}

Soit $S$ un schéma local nœthérien. Ratliff a démontré le théorème fondamental suivant, qui a déjà
été mentionné dans la proposition~\refext{I}{1Ratliff}.

\begin{theoreme2}[\cite{CRT@Matsumura} 31.7] \label{14-theoreme-Ratliff} Pour un schéma local
nœthérien $S$, les conditions suivantes sont équivalentes~: \begin{itemize} \item $S$ est
formellement caténaire, \item $S$ est universellement caténaire, \item $\Aff^1_S$ est caténaire,
\item $S$ est caténaire et pour tout $s \in S$, tout schéma intègre $S'$ muni d'une flèche finie et
dominante $S'\rightarrow \overline{\{s\}}$ et tout point fermé $s'$ de $S'$, on a $\dim(\OO_{S',s'})
= \dim(\overline{\{s\}})$.  \end{itemize} \end{theoreme2}

On a ajouté une quatrième condition équivalente à l'énoncé \cite[Theorem~31.7]{CRT@Matsumura}. Il
résulte de~\egalong{iv}{2}{5.6.10} que les trois premières conditions équivalentes impliquent la
quatrième. La réciproque est démontrée au cours de la démonstration de
\cite[Theorem~31.7]{CRT@Matsumura} (au second paragraphe de la page~255).


\begin{corollaire2}[\cite{CRT@Matsumura} 31.2] \label{14-corollaire-dim2-cat} Tout schéma nœthérien
de dimension~$\leq 2$ est caténaire. Tout schéma nœthérien de dimension~$\leq 1$ est universellement
caténaire.  \end{corollaire2}

\subsection{Schémas henséliens et caténarité}

Nous avons vu  que tout schéma local complet nœthérien est universellement caténaire dans la
proposition~\ref{14-prop-complet-univ-cat}. Les schémas locaux henséliens  jouissent également de
bonnes propriétés de caténarité.

\begin{proposition2} \label{14-prop-henselien-cat-univ-cat} Tout schéma local hensélien caténaire
est universellement caténaire.  \end{proposition2}

\begin{démo} Soit $S=\SP(A)$ un schéma local hensélien caténaire, soit $P$ un idéal premier de~$A$,
soit $L$ une extension finie de $\mathrm{Frac}(A/P)$ et soit $B$ une extension finie de $A/P$
contenue dans~$L$. D'après le théorème \ref{14-theoreme-Ratliff}, il suffit de prouver que la
dimension du localisé de $B$ en chacun de ses idéaux maximaux est égale à la dimension de $A/P$.
Toute algèbre finie sur un anneau hensélien est semi-locale d'après \egalong{iv}{4}{18.5} et
\egalong{iv}{4}{18.6}. Comme le schéma $B$ est intègre, il est local. Le théorème du « going-up »
(\cite{CRT@Matsumura}~9.3 et~9.4) montre qu'on a bien $\dim(B)=\dim(A/P)$.  \end{démo}

\medskip Rappelons également le résultat suivant, conséquence du
corollaire~\refext{I}{1-cor-normal-versus-complet}~ii).


\begin{proposition2} \label{14-prop-henselien-qexc-univ-cat} Tout schéma local hensélien
quasi-excellent est universellement caténaire.  \end{proposition2} Ainsi, tout schéma local
hensélien quasi-excellent est excellent.




\section{Spécialisations immédiates et fonctions de dimension}


\subsection{Définitions}

Soit $X$ un schéma. Pour tout point $x$ de $X$ et tout point géométrique $\bar{x}$ au-dessus de $x$,
on note $X_{(x)}$, $X_{(x)}^\h$ et $\widehat{X}_{(x)}$ le localisé, l'hensélisé et le complété de
$X$ en $x$. De même, l'on note $X_{(\bar{x})}$ l'hensélisé strict de $X$ en $\bar{x}$.

Soient $x$ et $y$ deux points de $X$, et $\bar{x}$ et $\bar{y}$ deux points géométriques au dessus
de $x$ et $y$. 

\begin{definition2}[\sga{4}{vii}{7.2}]\label{14-definition-specialisation} Un morphisme de
spécialisation $\bar{x} \spim\bar{y}$ est la donnée d'un $X$-morphisme $X_{(\bar x )} \rightarrow
X_{(\bar{y})}$ entre hensélisés stricts.  \end{definition2} D'après \sga{4}{vii}{7.4}, la donnée
d'une spécialisation  $\bar{x } \spim \bar{y}$ est équivalente à la donnée d'un $X$-morphisme
$\bar{x} \rightarrow X_{(\bar{y})}$.

\begin{definition2} \label{14-definition-specialisation-immediate} Soit $r \in \NN$. On dit qu'une
spécialisation $\bar{x} \spim \bar{y}$ est une  \emph{spécialisation de codimension $r$} si
l'adhérence de l'image de $\bar{x}$ dans $X_{(\bar{y})}$ est un schéma de dimension $r$.  \\ On dit
que $y$ est une \emph{spécialisation étale immédiate} de $x$ s'il existe une spécialisation $\bar{x}
\spim \bar{y}$ qui soit de codimension~$1$.  \\ On dit que $y$ est une \emph{spécialisation Zariski
immédiate} de $x$  si $y \in \overline{\{ x \} }$ et si le localisé en $y$ de l'adhérence de $x$
est de dimension $1$.  \end{definition2}

Si $y$ est une {spécialisation étale immédiate} de $x$, on dit également que $x$ est une
\emph{générisation étale immédiate} de $y$. Désignons par $f:X_{(\bar y)}\ra X_{(y)}$ le morphisme
d'hensélisation stricte. Les générisations  étales immédiates de $y$ sont alors les images par $f$
des points $x'\in X_{(\bar y)}$ tels que $\dim(\overline{\{x'\}})=1$.

Avant d'examiner plus en détail ces notions, on rappelle le fait facile suivant que nous utiliserons
implicitement plus bas : si $f\colon X\to S$ est un morphisme plat, $f$ envoie les points maximaux
de $X$ sur des points maximaux de $S$, autrement dit toute composante irréductible de $X$ domine une
composante irréductible de $S$.

Si $x$ et $y$ sont deux points d'un schéma nœthérien $X$ tels que $y\in \overline{\{x\}}$ (au sens
habituel, c'est dire que $y$ est une spécialisation de $x$ ou encore que $x$ est une générisation de
$y$), alors $y$ est une spécialisation Zariski (resp. étale) immédiate de $x$ si et seulement si
c'est le cas dans $\overline{\{x\}}_{(y)}$. Pour certaines considérations, ceci permet de supposer
que $X$ est local intègre de point générique $x$ et de point fermé $y$. Dans ce cas, $y$ est
spécialisation Zariski immédiate de $x$ si et seulement si $\dim(X)=1$. Dans le cas étale, cela se
lit sur l'hensélisé strict :

\begin{proposition2} \label{14-proposition-Laszlo-1} Si $x$ et $y$ sont deux points d'un schéma
nœthérien $X$, le point $y$ est une {spécialisation étale immédiate} de $x$ si et seulement si $y
\in \overline{\{ x \} }$ et que le hensélisé strict en un point géométrique au-dessus de $y$ de
l'adhérence de $x$ possède une composante irréductible de dimension $1$.  \end{proposition2}

\begin{démo} On se ramène au cas particulier $X=\overline{\{x\}}_{(y)}$ envisagé plus haut. Le point
$y$ est une spécialisation étale immédiate si et seulement s'il existe un point $\tilde{x}$ de
$X_{(\bar{y})}$ au-dessus de $x$ tel que l'adhérence de $\tilde{x}$ dans $X_{(\bar{y})}$ soit de
dimension $1$. Cela équivaut comme il est énoncé ici que $X_{(\bar{y})}$ possède une composante
irréductible de dimension $1$. En effet, si on note $\tilde{x}$ le point générique de $C$, par
l'argument de platitude énoncé plus haut, $C$ domine $X$, c'est-à-dire que $\tilde{x}$ est au-dessus
de $x$. Inversement, si $\tilde{x}$ est un point au-dessus de $x$ dont l'adhérence dans $X_{(\bar
y)}$ soit de dimension $1$, on peut noter $C$ une composante irréductible de $X_{(\bar y)}$
contenant $\tilde{x}$. Le point générique de $C$ et $\tilde{x}$ étant tous les deux au-dessus de
$x$, ils sont égaux puisque l'un est une générisation de l'autre et que les fibres de $X_{(\bar
y)}\to X$ sont discrètes.  \end{démo}

\medskip

\begin{proposition2} \label{14-prop-specialisation-zar-et} Soit $X$ un schéma nœthérien. Une
spécialisation Zariski immédiate entre points de $X$ est une spécialisation étale immédiate, et la
réciproque est vraie si $X$ est universellement caténaire.  \end{proposition2}

On peut supposer que $X=\overline{\{x\}}_{(y)}$ comme précédemment. Pour l'implication, on suppose
que $\dim(X)=1$ et on veut montrer que $X_{(\bar y)}$ possède une composante irréductible de
dimension $1$. En fait, toutes les composantes irréductibles de $X_{(\bar y)}$ sont de dimension
$1$. En effet, le schéma $X_{(\bar y)}$ est local et tout point de $X_{(\bar y)}$ autre que $\bar y$
appartient à la fibre de $x$ pour le morphisme \guil{pro-étale} $X_{(\bar y)}\to X$. L'espace
topologique sous-jacent à $X_{(\bar y)}-\{\bar y\}$ est donc fini, discret (et non vide). Ceci
implique que toute composante irréductible de $X_{(\bar y)}$ est de dimension $1$. Pour la
réciproque, nous utiliserons deux lemmes :

\begin{lemme2}\label{14-lemme-riou-1} Soit $X$ un schéma local nœthérien hensélien de point fermé
$y$. Soit $\bar y$ un point géométrique au-dessus de $y$. Alors, $X$ possède une composante
irréductible de dimension $1$ si et seulement si le hensélisé strict $X_{(\bar y)}$ en possède une.
\end{lemme2}

\begin{démo} Si $C$ est une composante irréductible de dimension $1$ de $X_{(\bar y)}$, son image
ensembliste dans $X$ est fermée car $p\colon X_{(\bar y)}\to X$ est entier. Comme $p$ est plat,
$p(C)$ est une composante irréductible de $X$ contenant exactement deux points donc $\dim(p(C))=1$.
Inversement, la surjectivité et la platitude de $p$ impliquent que si $D\subset X$ est une
composante irréductible de dimension $1$, il existe une composante irréductible $C$ de $X_{(\bar
y)}$ telle que $p(C)=D$. On a bien sûr $\dim(C)\geq 1$. Soit $z\in C$ un point qui ne soit pas le
point générique de $C$. Le point $p(z)$ ne peut pas être le point générique de $D$ car sinon la
fibre générique de $p$ ne serait pas discrète. C'est donc que $p(z)$ est le point fermé de $C$. Le
fait que $p^{-1}(y)$ soit discret implique alors que $z$ ne peut être que le point fermé de $C$.  Le
schéma local intègre $C$ possède donc exactement deux points : $\dim(C)=1$.  \end{démo}

\begin{lemme2}\label{14-lemme-riou-2} Soit $X$ un schéma local nœthérien. Si $X$ possède une
composante irréductible de dimension $1$, alors son complété $\widehat{X}$ aussi et la réciproque
est vraie si $X$ est universellement caténaire.  \end{lemme2}

\begin{démo} Commençons par le cas où $X$ est intègre. Comme $X$ et $\widehat{X}$ ont la même
dimension (cf. \cite[chap.~VIII, \S{}9]{CA@Zariski-Samuel}), si $\dim(X)=1$, alors
$\dim(\widehat{X})=1$. Un schéma local de dimension $1$ n'a bien sûr que des composantes
irréductibles de dimension $1$. Inversement, si $X$ est universellement caténaire, d'après le
théorème~\ref{14-theoreme-Ratliff}, les composantes irréductibles de $\widehat{X}$ ont toutes la
même dimension. Si l'une d'entre elles est de dimension $1$, le schéma $\widehat{X}$ est lui aussi
de dimension $1$, et alors $\dim(X)=\dim(\widehat{X})=1$.

Dans le cas général, notons $X_i$ les composantes irréductibles de $X$.  Pour tout $i$, le produit
fibré $X_i\times_X \widehat{X}$ s'identifie à $\widehat{X_i}$ (voir \sga{1}{iv}{3}). C'est un fait
que les composantes irréductibles des différents $\widehat{X_i}$ sont exactement les composantes
irréductibles de $\widehat{X}$ : ce sont des parties fermées irréductibles recouvrant $\widehat{X}$
et aucune d'entre elles n'est contenue dans une autre (ceci se déduit du fait que chaque composante
irréductible de $\widehat{X_i}$ domine $X_i$). Il est dès lors évident que l'énoncé pour $X$ résulte
de l'énoncé pour les schémas locaux intègres $X_i$.  \end{démo}

Montrons la réciproque énoncée dans la proposition~\ref{14-prop-specialisation-zar-et}. Soit $X$ un
schéma local nœthérien intègre universellement caténaire de point générique $x$ et de point fermé
$y$. On suppose que $y$ est une spécialisation étale immédiate de $x$. D'après la
proposition~\ref{14-proposition-Laszlo-1}, cela signifie que $X_{(\bar y)}$ possède une composante
irréductible de dimension $1$, ce qui équivaut d'après le lemme~\ref{14-lemme-riou-1} à dire que le
hensélisé $X_{(y)}^\h$ possède une composante irréductible de dimension $1$. Le complété
$\widehat{X}$ de $X$ est aussi celui de $X_{(y)}^\h$, le sens facile du lemme~\ref{14-lemme-riou-2}
appliqué à $X_{(y)}^\h$ montre que $\widehat{X}$ possède une composante irréductible de dimension
$1$. La réciproque de ce lemme appliquée au schéma universellement caténaire $X$ montre que $X$
possède une composante irréductible de dimension $1$. Ceci montre que $y$ est une spécialisation
Zariski immédiate de $x$.

\bigskip

Si le schéma $X$ est quasi-excellent, on peut alors lire les spécialisations étales d'un point $x$
de $X$ dans le complété de $X$ en $x$ :

\begin{proposition2}\label{14-prop-Laszlo-2} Soit $X$ un schéma quasi-excellent. Soient $x$ et $y$
deux points de $X$. On suppose que $y\in\overline{\{x\}}$. Notons $c:\widehat X_{(y)}\ra X_{(y)}$ le
morphisme de complétion. Le point $y$ est une spécialisation étale immédiate de $x$ si et seulement
si $c^{-1}(\overline{\{x\}})$ possède une composante irréductible de dimension $1$.
\end{proposition2}

\begin{démo} On peut supposer que $X=\overline{\{x\}}_{(y)}$. Le point $y$ est une spécialisation
étale immédiate de $X$ si et seulement si $X_{(\bar y)}$ possède une composante irréductible de
dimension $1$, c'est-à-dire, d'après le lemme~\ref{14-lemme-riou-1} que $X_{(y)}^\h$ en possède une.
Comme $X_{(y)}^\h$ est universellement caténaire (cf.
proposition~\ref{14-prop-henselien-qexc-univ-cat}), cela équivaut encore d'après le
lemme~\ref{14-lemme-riou-2} à dire que $\widehat{X}$ possède une composante irréductible de
dimension $1$.  \end{démo}

\medskip

\begin{definition2} \label{ddefi4} On appelle \emph{fonction de dimension} sur $X$ toute fonction
$\delta  : X \ra \ZZ $ telle que pour toute spécialisation étale immédiate $x \spim y$ entre points
de $X$, on ait $$\delta(y) = \delta(x) - 1.$$ \end{definition2}

La notion de fonction de dimension ne voit pas les éléments nilpotents~: $\delta$ est une fonction
de dimension sur~$X$ si et seulement si elle induit une fonction de dimension sur le sous-schéma
réduit~$X^{red}$. De plus si $U \stackrel{i}{\hookrightarrow } X$ est un morphisme étale et $\delta$
est une fonction de dimension sur $X$, $\delta\circ i$ définit une fonction de dimension sur $U$.
Plus précisément, l'ensemble des fonctions de dimension sur les $X$-schémas étales définit un
faisceau étale sur $X$. La différence entre deux fonctions de dimension sur $X$ est une fonction
invariante par spécialisations quelconques, donc une fonction localement constante. Nous montrerons
plus loin que si $X$ est quasi-excellent, des fonctions de dimension existent localement pour la
topologie étale sur $X$ et que les fonctions de dimension forment un $\ZZ$-torseur étale.

\subsection{Fonctions de dimension et universelle caténarité}

Le but de ce paragraphe est de démontrer le résultat suivant.

\begin{theoreme2}[Gabber]\label{14-theoreme-univ-cat-fonction-dim} Un schéma  nœthérien est
universellement caténaire si et seulement s'il possède une fonction de dimension localement pour la
topologie de Zariski.  \end{theoreme2}
 

Le théorème résulte de la conjonction du
corollaire~\ref{14-corollaire-univ-cat-fonction-dim-loc-zar} et de la
proposition~\ref{14-proposition-fonc-dim-univ-cat} ci-dessous. 

\begin{proposition2} \label{14-proposition--codim-fonc-dim} Soit $X$ un schéma intègre
universellement caténaire. La fonction $\delta:X\ra\ZZ$ définie  par $\delta(x)=-\dim(\OO_{X,x})$
est une fonction de dimension sur $X$.  \end{proposition2}

\begin{démo} En vertu de la proposition~\ref{14-prop-specialisation-zar-et}, comme $X$ est
universellement caténaire, il suffit de montrer que $\delta(y) = \delta(x) - 1$ pour toute
spécialisation Zariski immédiate~$x \spim y$. Comme $X$ est caténaire intègre, on a
$$\dim(\OO_{X,y}) = \dim(\OO_{X,x}) + \dim(\OO_{\overline{\{y\} },x}) = \dim(\OO_{X,x}) +1\: .$$
\end{démo}

\begin{remarque2} Si $X$ n'est pas supposé intègre, la fonction $\delta(x)=-\dim(\OO_{X,x})$ n'est
pas forcément une fonction de dimension comme le montre l'exemple où $X$ est obtenu par recollement
en un point d'une droite et d'un plan.  \end{remarque2}

\begin{corollaire2} \label{14-corollaire-univ-cat-fonction-dim-loc-zar} Tout schéma universellement
caténaire admet des fonctions de dimension localement pour la topologie de Zariski.
\end{corollaire2}


\begin{démo} Soit $X$ un schéma universellement caténaire. Soit $x\in X$. Il s'agit de montrer qu'il
existe un voisinage ouvert de $x$ pouvant être muni d'une fonction de dimension. L'espace
topologique $X$ est réunion de ses composantes irréductibles  $X_1,...,X_n$. Quitte à remplacer $X$
par l'ouvert complémentaire des  composantes $X_i$ ne contenant pas $x$, on peut supposer que $x$
appartient à toutes les composantes $X_i$.   Pour tout $0\leq i\leq n$,  notons $\mathscr{F}_i$
l'ensemble des fonctions de dimension sur $X_i$. D'après la
proposition~\ref{14-proposition--codim-fonc-dim}, cet ensemble est non vide et est un torseur sous
$\ZZ$. On choisit  un élément $\delta_i\in\mathscr{F}_i$ qui vaut $0$ sur le point $x$. Pour tous
$1\leq i , \; j \leq n$, la fonction $\delta_i-\delta_j$ est localement constante sur $X_i\cap X_j$
et vaut $0$ au point $x$. Soit $F_{i,j}$ le fermé de $X$, réunion des composantes connexes de
$X_i\cap X_j$ ne contenant pas $x$.  Soit $U$ le complémentaire dans $X$ de la réunion des
$F_{i,j}$.  Les fonctions $\delta_i$ se recollent en une fonction de dimension sur $U$.  \end{démo}

\medskip

Démontrons une réciproque  partielle du
corollaire~\ref{14-corollaire-univ-cat-fonction-dim-loc-zar}.

\begin{lemme2}\label{14-lemme-fonc-dim-cat} Un schéma nœthérien qui possède  une fonction de
dimension localement pour la topologie de Zariski est caténaire.  \end{lemme2}

\begin{démo} Pour montrer la caténarité, on peut supposer que le schéma $S$ possède une fonction de
dimension $\delta$. Supposons que $X\subset Y$ sont des fermés irréductibles de points génériques
respectifs $x$ et $y$. Choisissons une chaîne de spécialisations Zariski immédiates $y=x_0\spim
x_1\spim \dots x_d=x$ de longueur maximale. Par définition de la codimension, on a $\codim(X,Y)=d$
et par définition des fonctions de dimension, compte tenu du sens facile de la
proposition~\ref{14-prop-specialisation-zar-et}, on obtient $\delta(x)=\delta(y)-d$, d'où
$\codim(X,Y)=\delta(y)-\delta(x)$.

Maintenant, si $X\subset Y\subset Z$ sont des fermés irréductibles, on a :

\begin{eqnarray*} \delta(y)-\delta(x) &=& \codim(X,Y) , \\ \delta(z)-\delta(y) &=& \codim(Y,Z) , \\
\delta(z)-\delta(x) &=& \codim(X,Z) .  \end{eqnarray*}On en déduit
$\codim(X,Z)=\codim(X,Y)+\codim(Y,Z)$ d'où la caténarité.  \end{démo}

\medskip

Grâce au théorème~\ref{14-theoreme-Ratliff}, on peut remplacer \guil{caténaire} par
\guil{universellement caténaire} dans  le lemme~\ref{14-lemme-fonc-dim-cat} :

\begin{proposition2} \label{14-proposition-fonc-dim-univ-cat} Un schéma nœthérien qui possède une
fonction de dimension localement pour la topologie de Zariski est universellement caténaire.
\end{proposition2}

\begin{démo} On peut supposer qu'un tel schéma $S$ est local et possède une fonction de dimension
$\delta$. Notons $S^h$ son hensélisé. La fonction $\delta$ induit une fonction de
dimension~$\delta^h$ sur $S^h$. D'après le lemme \ref{14-lemme-fonc-dim-cat} appliqué à $S^h$ et à
$\delta^h$, le schéma $S^h$ est caténaire. D'après la proposition
\ref{14-prop-henselien-cat-univ-cat}, il est universellement caténaire. Le théorème
\ref{14-theoreme-Ratliff} montre alors que $S^h$ est formellement caténaire. Par conséquent, $S$ est
formellement caténaire donc également universellement caténaire.  \end{démo}

\subsection{Existence locale pour la topologie étale}

Dans ce paragraphe nous allons démontrer le théorème suivant.

\begin{theoreme2}[Gabber]\label{14-theoreme-exist-local-et-fn-dim} Tout schéma quasi-excellent
possède des fonctions de dimension localement pour la topologie étale.  \end{theoreme2}

Une application répétée du lemme suivant (variante de l'argument du
corollaire~\ref{14-corollaire-univ-cat-fonction-dim-loc-zar}) permet de montrer que si l'énoncé du
théorème est vrai pour les composantes irréductibles d'un schéma nœthérien $X$, alors le théorème
vaut aussi pour $X$. Plus loin, on pourra ainsi supposer que $X$ est intègre.

\begin{lemme2} Soit $X$ un schéma nœthérien dont l'espace topologique sous-jacent soit réunion de
deux sous-schémas fermés $X_1$ et $X_2$. Soit $\overline x$ un point géométrique de $X_1\cap X_2$.
On suppose que pour tout $i\in\{1,2\}$, il existe un voisinage étale $U_i$ de $\overline x$ dans
$X_i$ tel que $U_i$ admette une fonction de dimension.  Alors, il existe un voisinage étale $U$ de
$\overline x$ dans $X$ tel que $U$ admette une fonction de dimension.  \end{lemme2}

\begin{démo} Pour tout $i\in\{1,2\}$, on choisit un voisinage étale $U_i$ de $\overline x$ dans
$X_i$ tel que $U_i$ admette une fonction de dimension $\delta_i$. On se donne un point géométrique
distingué $\overline u_i$ au-dessus de $\overline x$ et on peut supposer que $\delta_i(u_i)=0$.
D'après \egalong{iv}{4}{18.1.1}, quitte à remplacer $U_i$ par un voisinage ouvert de $u_i$, on peut
supposer qu'il existe un morphisme étale $\widetilde{U_i}\to X$ et un isomorphisme
$\widetilde{U_i}\times_X{X_i}\simeq U_i$. On peut former le produit fibré $V=\widetilde{U_1}\times_X
\widetilde{U_2}$. Notons $\pi\colon V\to X$ la projection et $\overline{v}$ un point géométrique de
$V$ au-dessus de $\overline{u}_1$ et $\overline{u}_2$. Pour tout $i\in\{1,2\}$, la projection de $V$
sur le facteur $\widetilde{U_i}$ induit un morphisme étale $\pi^{-1}(X_i)\to U_i$. Par composition
avec ce morphisme étale, la fonction de dimension $\delta_i$ sur $U_i$ induit une fonction de
dimension $\tilde{\delta}_i$ sur le sous-schéma fermé $\pi^{-1}(X_i)$ de $V$ et elle vérifie
$\tilde{\delta}_i(v)=0$. Ces fonctions de dimensions $\tilde{\delta}_i$ pour $i\in\{1,2\}$ se
recollent sur l'ouvert $U$ complémentaire dans $V$ de la réunion des composantes connexes de
$\pi^{-1}(X_1\cap X_2)$ ne contenant pas $v$.  \end{démo}

Avant de traiter le cas des schémas intègres, commençons par celui des schémas normaux :


\begin{proposition2}\label{14-proposition-codim-fonction-de-dimension-sur-normal-quasi-excellent}
Soit $X$ un schéma normal quasi-excellent. La fonction $\delta : X \rightarrow \ZZ$ définie par
$\delta(x) = -\dim(\OO_{X,x})$ est une fonction de dimension.  \end{proposition2}


\begin{démo} On peut supposer de plus que $X$ est local. Notons $Y$ son hensélisé et $h\colon Y\to
X$ le morphisme de hensélisation.  D'après le théorème~\refext{I}{1th-greco} et les commentaires
subséquents, $Y$ est lui aussi quasi-excellent. D'après la
proposition~\ref{14-prop-henselien-qexc-univ-cat}, $Y$ est universellement caténaire. Par ailleurs,
comme le morphisme $Y\to X$ est régulier, la normalité de $X$ implique celle de $Y$ (cf.
\egalong{iv}{2}{6.5.4}). Le schéma local $Y$ est donc intègre et universellement caténaire, l'opposé
de la codimension définit une fonction de dimension $\delta'\colon Y\rightarrow \ZZ$. Comme une
spécialisation étale immédiate entre points de $X$ se relève pour ainsi dire par définition en une
spécialisation étale immédiate de points de $Y$, pour montrer que $\delta$ est une fonction de
dimension sur $X$, il suffit de montrer que pour tout $y\in Y$, si on note $x=h(y)$, on a
$\delta(x)=\delta'(y)$. Il s'agit donc de montrer que $\dim(Y_{(y)})=\dim(X_{(x)})$. Comme la
dimension d'un anneau local nœthérien ne change pas par complétion (et donc par hensélisation), il
suffit de montrer que $\dim(Y^\h_{(y)})=\dim(X^\h_{(x)})$. Les deux schémas locaux excellents
$Y^\h_{(y)}$ et $X^\h_{(x)}$ ayant le même hensélisé strict, ils ont bien la même dimension d'après
le lemme suivant : \end{démo}

\begin{lemme2} Soit $X$ un schéma local excellent. Notons $\overline x$ un point géométrique
au-dessus du point fermé $x$ de $X$. Alors, $\dim(X)=\dim X_{(\bar x)}$.  \end{lemme2}

\begin{démo} D'après le théorème~\ref{14-theoreme-univ-cat-fonction-dim}, il existe une fonction de
dimension $\delta$ sur $X$. On peut supposer qu'elle vérifie $\delta(x)=0$. Par composition avec le
morphisme $X_{(\bar x)}\to X$, on obtient une fonction de dimension $\delta'$ sur $X_{(\bar x)}$. En
utilisant le lien entre fonction de dimension et codimension énoncé dans la démonstration du
lemme~\ref{14-lemme-fonc-dim-cat}, il est évident que $\dim(X)$ (resp. $\dim(X_{(\bar x)})$) est le
maximum atteint par la fonction $\delta$ sur $X$ (resp. par $\delta'$ sur $X_{(\bar x)}$). Comme
$X_{(\bar x)}\to X$ est surjectif, le maximum de $\delta$ est le même que celui de $\delta'$, donc
$\dim(X)=\dim(X_{(\bar x)})$.  \end{démo}


\medskip

Revenons au cas du théorème~\ref{14-theoreme-exist-local-et-fn-dim} où $X$ est supposé intègre et
quasi-excellent, et notons $Y$ son normalisé. Le morphisme $p : Y \rightarrow  X$ est fini et
surjectif, donc de descente cohomologique universelle. Notons~$\delta$ une fonction de dimension sur
$Y$ ; son existence est assurée par la
proposition~\ref{14-proposition-codim-fonction-de-dimension-sur-normal-quasi-excellent}. Soit~$p_1$
et~$p_2$ les deux projections $Y\times_X Y \ra Y$. La fonction $$p_1^\star \delta -p_2^\star\delta :
Y\times_X Y \longrightarrow \ZZ$$ définit un $1$-cocycle de Čech, donc une classe~$[p_1^\star
\delta -p_2^\star\delta]$ dans $\HH^1_{\v{C}ech}(Y \ra X, \ZZ)$.  D'après la théorie de la descente
cohomologique, il existe une injection naturelle $$\HH^1_{\v{C}ech}(Y \ra X, \ZZ) \hookrightarrow
\HH^1(X,\ZZ)\:.$$ La classe~$[p_1^\star \delta -p_2^\star\delta]$ définit donc une classe
d'isomorphisme de $\ZZ$-torseurs étales sur~$X$. Il résulte alors immédiatement de la proposition
suivante que $X$ admet une fonction de dimension localement pour la topologie étale :

\begin{proposition2} Soit $U$ un schéma étale sur $X$. L'annulation de la classe $[p_1^\star \delta
-p_2^\star\delta]\vert_U$ dans $\HH^1(U,\ZZ)$ entraîne l'existence d'une fonction de dimension sur
$U$.  \end{proposition2}

\begin{démo} En utilisant la compatibilité des constructions au changement de base étale $U\to X$,
on peut supposer que $U=X$.  L'annulation de $[p_1^\star \delta -p_2^\star\delta]$ dans
$\HH^1_{\v{C}ech}(Y \ra X,\ZZ) \hookrightarrow \HH^1(X,\ZZ)$ signifie qu'il existe une fonction
\emph{localement constante} $\gamma: Y \ra \ZZ$ telle que $p_1^\star \delta-p_2^\star
\delta=p_1^\star \gamma-p_2^\star \gamma$. Autrement dit, quitte à remplacer $\delta$ par
$\delta-\gamma$, on peut supposer que $p_1^\star \delta=p_2^\star \delta$. Ainsi, $\delta:Y\ra \ZZ$
se descend en une fonction $\delta':X\ra \ZZ $.

Pour conclure, il s'agit de montrer que si $p: Y\ra X$ est un morphisme fini surjectif entre schémas
quasi-excellents, que $\delta'\colon X\to \ZZ$ est une fonction et $\delta=\delta'\circ p$, alors
$\delta'$ est une fonction de dimension sur $X$ si $\delta$ est une fonction de dimension sur $Y$.
Pour cet énoncé, on peut supposer que $X$ est local hensélien quasi-excellent, donc universellement
caténaire (cf.  proposition~\ref{14-prop-henselien-qexc-univ-cat}), ce qui permet d'appliquer la
proposition~\ref{14-prop-specialisation-zar-et}. Pour montrer que $\delta'$ est une fonction de
dimension sur $X$, il suffit alors de savoir qu'au-dessus d'une spécialisation Zariski immédiate
$s\spim t$ de points de $X$, il existe une spécialisation Zariski immédiate $s'\spim t'$ de points
de $Y$.  \end{démo}

\subsection{Existence globale de fonctions de dimension}


Suivant \egazero{iv}{14.2.1}, on dit qu'un schéma nœthérien $X$ est équicodimensionnel si ses points
fermés ont tous la même codimension (qui est alors égale à $\dim(X)$).

\begin{exemple2} Les schémas de type fini équidimensionnels sur un corps $k$ ou sur $\ZZ$ sont
équicodimensionnels~: il est classique que dans cette situation, on a $\dim(X) = \dim(\OO_{X,x})$
pour tout point fermé $x$.  Les schémas locaux sont équicodimensionnels car ils possèdent un unique
point fermé.  Si $S=\mathrm{\SP(R)}$ est un trait d'uniformisante $\pi$, le schéma $\Aff^1_S$ n'est
pas équicodimensionnel. En effet, il existe un point fermé de $\Aff^1_S$ au-dessus du point
générique de $S$~:~il suffit d'écrire $\Aff^1_S=\SP(R[t])$ et de considérer $\MM=(\pi t -1 )$, qui
est un idéal maximal de corps résiduel $\mathrm{Frac}(R)$.  \end{exemple2}

Le lemme suivant est inspiré de~\egazero{iv}{14.3.3}\;\footnote{Gabber remarque que la proposition
\egazero{iv}{14.3.3} est fausse. Les assertions $a$, $c$ et $d$ de \emph{loc. cit.} sont
équivalentes entre elles et impliquent $b$ mais ne lui sont pas équivalentes. Il faut remplacer $b$
par la condition « $X$ est caténaire équidimensionnel et ses composantes irréductibles sont
équicodimensionnelles ». Gabber donne comme contre-exemple le spectre du localisé de
$k[x,y,z,w]/(xz,xw)$ en le complémentaire de l'union des idéaux premiers $(x-1,y)$ et $(x,z,w)$ avec
$k$ un corps.}.

\begin{lemme2} \label{14-lemme-EGA-faux} Soit $X$ un schéma équidimensionnel caténaire dont les
composantes irréductibles sont équicodimensionnelles. Pour tout $x\in X$, on a $$\dim(X) =
\dim\left(\overline{\{x\}}\right)+\dim(\OO_{X,x})\: .$$ \end{lemme2}

\begin{remarque2} En particulier, cette égalité est vérifiée pour tout schéma intègre local
caténaire. D'après~\cite[th.~31.4]{CRT@Matsumura}, si $X$ est intègre local nœthérien et si pour
tout $x\in X$, on a  $\dim(X)=\dim(\overline{\{x\}} )+\dim(\OO_{X,x})$, alors $X$ est caténaire.
\end{remarque2}

Le lemme~\ref{14-lemme-EGA-faux} et la proposition~\ref{14-proposition--codim-fonc-dim} impliquent
le résultat suivant.

\begin{corollaire2} \label{14-coro-equicodim-fonction-dim} Soit $X$ un schéma intègre,
équicodimensionnel et universellement caténaire. La fonction $\delta:X\ra\ZZ$ définie par $\delta(x)
= \dim(\overline{\{x\}})$ est une fonction de dimension sur $X$.  \end{corollaire2}

Les conclusions du corollaire sont prises en défaut si $X$ n'est pas équicodimensionnel. Soient par
exemple $S=\SP(R)$ un trait d'uniformisante $\pi$ et $X=\Aff^1_S=\SP(R[t])$. Si l'on note $x$ le
point fermé de $X$ correspondant à l'idéal maximal $(\pi t -1)$ et  $\eta$  le point générique de
$\Aff^1_S$, alors la spécialisation $\eta\spim x$ est immédiate et pourtant
$\dim(\overline{\{x\}})=0$ et $\dim(\overline{\{\eta\}})=2$.

\begin{corollaire2}\label{14-coro-schema-tf-sur-corps-fonc-dim} Soit $X$ un schéma qui est soit de
type fini sur un corps, soit de type fini sur $\ZZ$, ou soit local universellement caténaire. La
fonction définie par $\delta(x) = \dim(\overline{\{x\}})$ est une fonction de dimension sur $X$.
\end{corollaire2}

\begin{démo} Le schéma $X$ est universellement caténaire. D'après le
corollaire~\ref{14-coro-equicodim-fonction-dim}, la fonction $\delta$ est une fonction de dimension
sur chaque composante irréductible de $X$. Cette fonction est définie globalement donc est une
fonction de dimension sur $X$.  \end{démo}


\subsection{Fonction de dimension induite}

Soient $Y\ra X$ un morphisme de schémas et $\delta_X$ une fonction de dimension sur $X$. Dans
certains cas nous  pouvons construire une fonction de dimension~$\delta_Y$ induite sur $Y$. On admet
la proposition suivante.

\begin{proposition2}[\cite{CA@Matsumura},~14.C]\label{14-Matsumura-14-C} Soient $X$ un schéma
nœthérien intègre universellement caténaire, $Y$ un schéma intègre et $Y\rightarrow X$ un morphisme
de type fini dominant. Soient $k(X)$ et $k(Y)$ les corps de fractions respectifs de $X$ et $Y$,
soient $y$ un point de $Y$ et $x$ son image dans $X$, et soient $k(y)$ et $k(x)$ leurs corps
résiduels. On a $$\dim(\OO_{Y, y}) - \mathrm{degtr}\left(k(Y)/k(X)\right)  =   \dim(\OO_{X, x})-
\mathrm{degtr}\left(k(y)/k(x)\right) \: .$$ \end{proposition2}

\begin{corollaire2} \label{14-coro-fonction-dim-induite-tf} Soient $X$ un schéma nœthérien qui
possède une fonction de dimension $\delta_X$ et $f  :  Y \ra X$ un morphisme de type fini. La
fonction $\delta_Y  :  Y \ra \ZZ$ définie par $$\delta_Y(y) =  \delta_X(f(y))  +
\mathrm{degtr}\left(k(f(y))/k(y)\right)$$ est une fonction de dimension sur $Y$.  \end{corollaire2}


\begin{démo} On peut supposer que $X$ et $Y$ sont intègres et que $f$ est dominant. D'après la
proposition \ref{14-proposition-fonc-dim-univ-cat}, $X$ est universellement caténaire et d'après la
proposition  \ref{14-proposition--codim-fonc-dim}, $x\mapsto -\dim(\OO_{X,x})$ est une fonction de
dimension sur $X$. Comme les fonctions de dimension forment un $\ZZ$-torseur, on peut supposer que
$\delta_X(x)  =  -\dim(\OO_{X,x})$ pour tout $x\in X$.

Le corollaire~\ref{14-Matsumura-14-C} montre que $\delta_Y(y) =-\dim(\OO_{Y,y}) +
\mathrm{degtr}(k(Y)/k(X))$ et la proposition~\ref{14-proposition--codim-fonc-dim} montre que $y
\mapsto -\dim(\OO_{Y,y})$ est une fonction de dimension sur $Y$. Ainsi, $\delta_Y$ est une fonction
de dimension sur $Y$.  \end{démo}

\medskip Avant d'établir la fonctorialité des fonctions de dimension vis-à-vis des morphismes
réguliers entre schémas excellents, démontrons un énoncé de changement de base par un morphisme
régulier en cohomologie étale. Ce lemme est une simple conséquence du théorème de Popescu
\refext{I}{1-Artin-Popescu} et du théorème de changement de base par un morphisme
lisse~\sga{4}{xvi}{1.2}.

\begin{lemme2} \label{14-lemme-cb-regulier} Soient $$\xymatrix{T' \ar[r]^{g'} \ar[d]_{f'} & T
\ar[d]^f \\ S' \ar[r]^g & S}$$ un diagramme cartésien de schémas, $n$ un entier inversible sur~$S$
et~$\mc{F}$ un faisceau étale en~$\ZZ/n\ZZ$-modules sur~$T$. Supposons que $f$ est quasi-compact et
quasi-séparé et que $g$ est un morphisme régulier entre schémas nœthériens. La flèche naturelle de
changement de base $$g^* \R f_*  (\mc{F})  \isolong   \R f'_* g'^*  (\mc{F})$$ est un isomorphisme.
\end{lemme2}

\begin{démo} D'après le théorème de Popescu, il existe un ensemble ordonné filtrant~$I$ et une
famille de schémas~$S_i$ indexée par~$I$, tels que~$S_i$ soit lisse sur~$S$ pour tout~$i\in I$ et
que $S' = \lim_i \: S_i$. Il existe donc pour tout~$i\in I$ un diagramme commutatif à carrés
cartésiens $$\xymatrix{T' \ar[r]^{h'_i} \ar[d]^{f'} & T_i \ar[r]^{g'_i} \ar[d]^{f_i} & T \ar[d]^f \\
S' \ar[r]^{h_i} & S_i \ar[r]^{g_i} & S}$$ On conclut grâce à la suite d'isomorphismes suivante pour
tout $q\geq 0$ : \begin{eqnarray*} \R^q f'_*  g'^*  (\mc{F})  & \isolongg & {\colim}_{i} \: h_i^*
\R^q {f_i}_*  g_i'^* (\mc{F}) \\ & \isolongg &  {\colim}_{i} \: g^*  \R^q f_* (\mc{F}) \\ &
\isolongg & g^*  \R^q f_* (\mc{F}) \end{eqnarray*} Le premier de ces isomorphismes résulte du
théorème de passage à la limite~\sga{4}{vii}{5.11}, et le second du théorème de changement de base
par le morphisme lisse~$g_i$~\sga{4}{xvi}{1.2}.  \end{démo}

Nous prouvons à présent qu'un morphisme régulier entre schémas excellents permet d'induire des
fonctions de dimension.

\begin{proposition2} \label{14-prop-fonction-dim-induite-regulier} Soient $f : Y \ra X$ un morphisme
régulier entre schémas excellents et~$\delta_X$ une fonction de dimension sur~$X$. La fonction
$\delta_Y : Y \ra \ZZ$ définie par $$\delta_Y(y) =  \delta_X(f(y))  - \dim \left( \OO_{Y_{f(y)},y}
\right) $$ est une fonction de dimension sur $Y$.  \end{proposition2}


\begin{démo} Comme la vérification est locale, il n'y a pas de mal à supposer $X$ et $Y$ strictement
locaux et $f$ local. Soit $\delta$ une fonction de dimension sur $Y$ ; son existence est assurée par
le théorème \ref{14-theoreme-univ-cat-fonction-dim}. Il suffit de montrer que $\delta_Y-\delta$ est
une fonction constante sur $Y$. Les fibres de~$f$ sont  régulières donc universellement caténaires
d'après~\ref{14-exemple-regulier-implique-univ-cat}. La
proposition~\ref{14-proposition--codim-fonc-dim} montre que la fonction qui à~$y\in Y$ associe  $$-
\dim \left( \OO_{Y_{f(y)},\: y} \right) $$ induit une fonction de dimension sur chacune des fibres
de~$f$. La fonction $\delta_Y-\delta$ est donc localement constante sur chaque fibre de~$f$. Il
résulte du lemme~\ref{14-lemme-cb-regulier} que ces fibres sont connexes : en effet, on a
$\HH^0(f^{-1}(x),\ZZ/n\ZZ)=\HH^0(x,\ZZ/n\ZZ)=\ZZ/n\ZZ$ pour tout $x\in X$ et tout entier~$n$
inversible sur~$X$. La fonction $\delta_Y-\delta$ est donc constante sur les fibres de~$f$ et
descend à $X$. Il suffit de montrer que~$\gamma=\delta_Y-\delta$ est localement constante sur $X$.
Une façon de calculer la valeur de $\gamma$ en un point $s$ de $X$ consiste à considérer le point
générique $\eta_s$ du schéma régulier connexe $f^{-1}(s)$, de sorte que
$\gamma(s)=\delta_X(s)-\delta(\eta_s)$. Soit $s' \spim s$ une spécialisation Zariski immédiate entre
deux points de $X$. Il s'agit de montrer que $\gamma(s)=\gamma(s')$. Vu que $\delta_X$ et $\delta$
sont des fonctions de dimension sur $X$ et $Y$ respectivement, pour montrer cela, il suffit de
savoir que $\eta_{s'}$ est une spécialisation immédiate de $\eta_s$. Pour montrer cela, quitte à
remplacer $X$ par le localisé en $s$ de l'adhérence de $s'$, on peut supposer que $X$ est local
intègre de dimension $1$, de point générique $s'$ et de point fermé $s$. Il s'agit alors de montrer
que la fibre $f^{-1}(s)$ est de codimension $1$ dans $Y$, ce qui résulte facilement du
\emph{Hauptidealsatz}.  \end{démo}

\subsection{Contre-exemple}

Rappelons l'exemple de~\ega{ii}{5.6.11} d'un schéma caténaire non universellement caténaire.  Soient
$k_0$ un corps et $k$ une extension purement transcendante de $k_0$  de degré de transcendance
infini. Notons $S=k[X]_{(X)}$ le localisé de l'anneau de la droite affine sur~$k$ en l'origine et
$V=S[T]$. Les idéaux maximaux $\MM=(X,T)$ et $\MM'=(XT-1)$ de $V$ sont respectivement de hauteur $2$
et $1$, et il existe un isomorphisme $\phi \: : \: V/\mathfrak{m} \iso V/\mathfrak{m}'$. Posons
$C=\{\:v\in V \: | \: \phi(\: v \: \mathrm{mod} \: \mathfrak{m}'\:) = v \: \mathrm{mod} \:
\mathfrak{m} \:\}$. C'est un sous-anneau de $V$ qui n'est pas de type fini sur $k$ et dont le
spectre est le quotient de $\SP(V)$ par la relation d'équivalence identifiant $\MM$ et $\MM'$. Il
est intègre et contient l'idéal maximal $\mathfrak{n}$ image réciproque de $\MM$.

\begin{proposition2}\label{14-propo-contre-ex} Le schéma $\SP(C)$ est nœthérien, quasi-excellent,
caténaire mais  non universellement caténaire. Le point fermé correspondant à l'idéal maximal
$\mathfrak{n}$ de $C$ est une spécialisation étale immédiate mais non Zariski immédiate du point
générique de $\SP(C)$.  \end{proposition2}

\begin{démo} Le caractère nœthérien est montré dans \ega{ii}{5.6.11} et le caractère quasi-excellent
dans \ega{ii}{7.8.4}. Le schéma $\SP(C)$ est caténaire d'après le
corollaire~\ref{14-corollaire-dim2-cat} car il est de dimension~$2$. On montre qu'il n'est pas
universellement caténaire en vérifiant que la dernière condition du
théorème~\ref{14-theoreme-Ratliff} est prise en défaut. En effet, le morphisme
$$\SP(V_{(\mathfrak{n})}) \ra \SP(C_{(\mathfrak{n})})$$ est fini  et la fibre du point fermé
correspondant à $\mathfrak{n}$ est constituée des points fermés correspondants à $\mathfrak{m}$ et
$\mathfrak{m}'$, qui sont de hauteur respective $1$ et $2$. Le schéma $\SP(C)$ est intègre de
dimension $2$ donc $\mathfrak{n}$ n'est pas une spécialisation Zariski immédiate du point générique.
\end{démo}

\ifx\danslelivre\undefined
\bibliography{../fichiers-configuration-latex/bibliographie}
\bibliographystyle{../fichiers-configuration-latex/smfalpha}
\end{document}
\fi
